\documentclass[12pt]{article}
\usepackage{amsmath,amssymb}
\begin{document}
\section*{Research notes}

Problem: uniqueness of invariant probability for Dirichlet stochastic heat
equation on $[0,1]$ with additive space-time white noise and distributional
drift $\delta_0(u)$, under the ABLM regularized/mild interpretation and
regularisation by $G_{1/n}$.

Key interpretation:
\[
\delta_0(r)=-\frac12(-2{\bf1}_{r>0})',
\]
so the equation is a singular gradient SPDE with bounded potential
$f(r)=-2H_+(r)$.  For smooth approximation $B_n'=G_{1/n}$ with
$B_n(r)=\int_{-\infty}^rG_{1/n}(s)\,ds$, the invariant measure of the
regularized gradient equation is
\[
\pi_n(dv)=Z_n^{-1}\exp\left(2\int_0^1 B_n(v(x))\,dx\right)\gamma(dv),
\]
where $\gamma$ is Brownian bridge law (the invariant measure of the
Dirichlet OU/string process).  Since $0\le B_n\le1$ and $B_n\to H_+$
away from zero, while a Brownian bridge spends zero Lebesgue time at
zero,
\[
\pi_n\to
\pi(dv)=Z^{-1}\exp\left(2\int_0^1{\bf1}_{v(x)>0}\,dx\right)\gamma(dv)
\]
in total variation.  Densities $d\pi_n/d\gamma$ and $d\pi/d\gamma$ are
bounded above and below by deterministic positive constants.

Round-1 robust result obtained: assuming pointwise weak convergence of
regularized semigroups on bounded continuous observables,
\[
P_t^n\Phi(x)\to P_t\Phi(x),\qquad x\in C_0,\ t>0,
\]
one can pass reversibility/invariance and a uniform spectral gap from
$P^n$ to the actual ABLM semigroup $P$ without invoking
Bounebache--Zambotti:
\[
\int gP_tf\,d\pi=\int fP_tg\,d\pi,\qquad
\operatorname{Var}_\pi(P_tf)\le e^{-2ct}\operatorname{Var}_\pi(f).
\]
Key estimates:
\[
\operatorname{Var}_{\pi_n}(F)\le C\mathcal E_n(F,F),\qquad
\mathcal E_n(F,F)=\frac12\int\|DF\|_{L^2}^2\,d\pi_n,
\]
with $C$ uniform in $n$, by Gaussian Poincare for $\gamma$ plus bounded
density perturbation.  Then pass the $L^2$ contraction/spectral gap to
the limit using total variation convergence $\pi_n\to\pi$ and bounded
pointwise convergence $P_t^nf\to P_tf$.

Abstract uniqueness mechanism:
If $\eta$ is invariant and the actual resolvent satisfies
\[
R_\lambda^P(x,\cdot)\ll\pi\quad\hbox{for every }x,
\]
then $\eta=\lambda\eta R_\lambda^P\ll\pi$.  With $h=d\eta/d\pi$,
symmetry gives $P_th=h$ in $L^1(\pi)$.  For $h_a=h\wedge a$, concavity
gives $h_a\ge P_th_a$ and equal integrals force equality.  Since
$h_a\in L^2$, the spectral gap forces $h_a$ constant.  Varying $a$
implies $h$ constant, hence $\eta=\pi$.

Central unresolved issue:
Need pointwise resolvent absolute continuity for the actual ABLM
semigroup, not only for the $L^2(\pi)$ Dirichlet-form version and not
only quasi-everywhere.  BZ Proposition 2.5 gives an $L^2$ resolvent kernel
$\rho_\lambda(x,dy)\ll\nu(dy)$ for all $x$ (setting it to zero on an
exceptional set) and transition absolute continuity for all
non-exceptional initial conditions.  BZ Proposition 3.1 constructs a
weak local-time solution for quasi-every initial condition.  ABLM
uniqueness plus the 2025 weak/mild equivalence paper plausibly identify
the BZ process with the ABLM mild process for quasi-every initial point,
but this still does not prove that arbitrary invariant measures cannot
charge the exceptional set, nor that $R_\lambda^P(x,\cdot)\ll\pi$ for
$x$ in that set.

Potential routes to close the gap:
1. Uniform heat-kernel/capacity estimates for $P_t^n$: prove
   $P_t^n(x,A)\le C_{t,x}\,\mathrm{Cap}_n(A)^\alpha$ or
   $\|dP_t^n(x,\cdot)/d\pi_n\|_{L^2(\pi_n)}\le C_{t,x}$ with constants
   stable as $n\to\infty$ (perhaps only for $t$ large).  This would pass
   AC to the limit and imply uniqueness via convergence to equilibrium.
2. Strict version of BZ process: prove that the ABLM process is the
   strict all-starting-points Hunt process associated with the BZ
   Dirichlet form, and that BZ's resolvent kernel represents its
   resolvent at every deterministic $x$.
3. Entrance from exceptional sets: let $N$ be the BZ exceptional set.
   Show $P_s^P(x,N)=0$ for every $x$ and $s>0$.  Then q.e. resolvent AC
   plus the Markov property gives all-$x$ resolvent AC:
   $R_\lambda(x,A)\le\int_0^s e^{-\lambda t}P_t(x,A)dt$ for
   $\pi(A)=0$, then let $s\downarrow0$.
4. Prove uniqueness of invariant measures without pointwise AC by showing
   every invariant stationary solution has enough finite-energy/Sobolev
   regularity not to charge zero-capacity sets.  Stationarity and the
   nondegenerate additive noise imply one-dimensional projections have no
   atoms by the occupation formula, but this is far weaker than avoiding
   all Gaussian-capacity-zero sets.

Useful literature facts:
- BZ use $\nu(dx)\propto \exp(-\int f(x_r)dr)\mu(dx)$.  For
  $f=-2H_+$ this is exactly $\pi$.
- BZ equation has singular term
  $-\frac12\int f(da)\partial_\theta\ell^a_{t,\theta}$; for
  $f=-2H_+$ this is $\partial_\theta\ell^0$, formally
  $\delta_0(u(t,\theta))$.
- ABLM 2024 proves strong existence/uniqueness for $b=\kappa\delta_0$
  and convergence of smooth approximations; it explicitly did not treat
  Dirichlet boundary conditions, but the problem assumes the extension.
- ABLM 2025 ``Analytically weak and mild solutions...'' proves weak
  regularized and mild regularized solution notions coincide in the
  known existence range; may be useful for local-time-to-mild
  identification, though it still does not by itself remove the
  exceptional-set issue.



\section*{Additional reduction lemmas and strategy notes}

Strong-Feller/resolvent-Feller reduction.  Since $\pi$ is invariant and
has full support, if $A$ is Borel and $\pi(A)=0$, then
\[
  \int P_t{\bf1}_A\,d\pi=\pi(A)=0,\qquad
  \int R_\lambda{\bf1}_A\,d\pi=\lambda^{-1}\pi(A)=0.
\]
Thus either of the following immediately upgrades $\pi$-a.e. null-set
avoidance to all starting points:
\[
  P_t:B_b(E)\to C_b(E)\quad \hbox{for every }t>0,
\]
or
\[
  R_\lambda:B_b(E)\to C_b(E)\quad \hbox{for some }\lambda>0.
\]
Indeed a nonnegative continuous function with zero integral against a
full-support probability is identically zero.  This proves
$P_t(x,\cdot)\ll\pi$ for all $x,t>0$ in the strong-Feller case, and
$R_\lambda(x,\cdot)\ll\pi$ in the resolvent-Feller case.

Entrance reduction.  Let $N$ be a BZ exceptional set such that the ABLM
and BZ resolvents are identified for $y\notin N$.  If
\[
  P_s^P(x,N)=0,\qquad x\in E,\ s>0,
\]
then all-starting-point resolvent absolute continuity follows.  For
$\pi(A)=0$,
\[
 R_\lambda(x,A)
 =\int_0^\varepsilon e^{-\lambda t}P_t(x,A)\,dt
  +e^{-\lambda\varepsilon}P_\varepsilon R_\lambda(\cdot,A)(x).
\]
The first term is at most $\varepsilon$ and the second term is at most
$\lambda^{-1}P_\varepsilon(x,N)$, hence zero; let
$\varepsilon\downarrow0$.

ABLM method note.  The ABLM 2024 paper explicitly says that its proof
uses neither Girsanov nor Zvonkin transformations; therefore one should
not rely on an unstated global Zvonkin homeomorphism to transfer strong
Feller or transition equivalence from an OU process.  Any strong-Feller
argument would have to be proved separately, probably using Malliavin,
capacity, or stochastic-sewing/Krylov estimates.

OU comparison.  For the linear Dirichlet stochastic heat equation,
\[
  Z_t=S_t x+\int_0^t S_{t-s}\,dW_s,
\]
the law at time $t>0$ is equivalent to $\gamma$.  In the sine basis with
Dirichlet eigenvalues $\lambda_k=(k\pi)^2$, the covariance eigenvalues are
$q_k(t)=(1-e^{-\lambda_k t})/\lambda_k$ up to the convention-dependent
factor, while $\gamma$ has eigenvalues proportional to $\lambda_k^{-1}$.
The covariance ratio differs from $1$ by a summable sequence and the mean
$S_t x$ is in the Cameron--Martin space $H^1_0$, so Feldman--Hajek gives
equivalence.  Thus it is enough, in principle, to control the singular
drift contribution relative to this Gaussian law.

Uniform-density route.  A sufficient estimate for the regularized
semigroups is: for every $t>0$ and $x\in E$, there exist $p>1$ and
$C_{t,x}<\infty$ such that
\[
  \sup_n \left\|
    \frac{dP_t^n(x,\cdot)}{d\pi_n}
  \right\|_{L^p(\pi_n)}\le C_{t,x}.
\]
Together with $\pi_n\to\pi$ in total variation and weak convergence
$P_t^n(x,\cdot)\Rightarrow P_t(x,\cdot)$, this implies
$P_t(x,\cdot)\ll\pi$.  A capacity version
\[
  P_t^n(x,O)\le C_{t,x}\operatorname{Cap}_n(O)^\alpha
\]
uniformly in $n$ would also imply entrance from BZ exceptional sets,
because the forms have densities bounded above and below relative to
Brownian bridge and hence comparable capacities.

Naive Girsanov issue.  For fixed $n$, the regularized drift
$G_{1/n}(u)$ is bounded and $H$-valued, so pathwise Girsanov from the
linear equation gives absolute continuity.  The usual Novikov/entropy
term contains
\[
  \int_0^t\|G_{1/n}(u_s)\|_{L^2(0,1)}^2\,ds,
\]
which may blow as $n\to\infty$ because $G_{1/n}^2$ approximates a
multiple of $\sqrt n\,\delta_0$.  Any successful limiting argument must
use cancellation or estimates depending on the bounded potential
$0\le B_n\le1$, not on $\|G_{1/n}\|_\infty$ or $\|G_{1/n}\|_2$.

Dirichlet-form comparison idea.  The regularized and limiting invariant
densities are uniformly bounded above and below relative to Brownian
bridge.  Hence the forms
\[
  \mathcal E_n(F,F)=\frac12\int\|DF\|_{L^2}^2\,d\pi_n,\qquad
  \mathcal E_0(F,F)=\frac12\int\|DF\|_{L^2}^2\,d\gamma
\]
have comparable $1$-energies and capacities.  This gives spectral-gap
and compact-resolvent information uniformly, but by itself only produces
kernels and absolute continuity in the $L^2(\pi)$ or quasi-everywhere
sense.  The missing upgrade is a strict, all-starting-points statement.



\section*{Additional reductions for the exceptional-set gap}

Eventual equicontinuity/e-property route.  Suppose $P_t$ is already known
to be reversible with invariant $\pi$, and $P_t f\to\pi(f)$ in
$L^2(\pi)$ for every bounded Lipschitz $f$.  If for every $x\in E$ and
$\varepsilon>0$ there are a neighbourhood $U$ of $x$ and $t_0$ such that
\[
  |P_t f(y)-P_t f(x)|\le \varepsilon,\qquad y\in U,\quad t\ge t_0,
\]
then $P_t f(x)\to\pi(f)$ for every $x$.  The proof avoids compactness:
since $\pi$ has full support, $\pi(U)>0$, and
\[
 |P_t f(x)-\pi(f)|
 \le \varepsilon+
 \left(\pi(U)^{-1}\int_U |P_t f(y)-\pi(f)|^2\,\pi(dy)\right)^{1/2}.
\]
The second term goes to $0$.  Then any invariant law $\eta$ satisfies
$\eta(f)=\int P_t f\,d\eta\to\pi(f)$; bounded Lipschitz functions
separate probability measures on $E$.  This route would prove uniqueness
without all-starting-point absolute continuity.  The analytic missing
input is the eventual equicontinuity estimate.  Ordinary finite-time
Feller continuity is not enough, because the modulus may deteriorate as
$t\to\infty$.

Possible use of ABLM stability estimates.  In the ABLM uniqueness proof
for two solutions driven by the same noise and with the same initial
condition, the local estimate has the form
\[
\|z_t-P_{t-s}z_s\|\lesssim
D_{s,t}(t-s)^{1/2+\delta}+D_{s,t}|\log D_{s,t}|(t-s),
\]
where $D_{s,t}=\sup_{r\in[s,t]}\|z_r\|$.  For different initial data the
same local identity should still hold for the nonlinear drift part,
while $P_tz_0$ remains as an inhomogeneous term.  This suggests finite
time continuous dependence (hence weak Feller), perhaps by an Osgood
argument.  It does not immediately imply the eventual equicontinuity
needed above; one would need either a uniform-in-time modulus or an
asymptotic-strong-Feller estimate in which the heat semigroup
dissipation beats the logarithmic loss.

Capacity-smooth invariant measures route.  Let $N$ be the BZ exceptional
set on whose complement the resolvent kernel represents the actual
resolvent and is absolutely continuous with respect to $\pi$.  It is
enough to show that every invariant probability $\eta$ is smooth for the
BZ Dirichlet form, in particular $\eta(N)=0$.  A sufficient finite-energy
inequality is
\[
  \left|\int \varphi\,d\eta\right|
  \le C_\eta \mathcal E_1(\varphi,\varphi)^{1/2}
\]
for bounded nonnegative quasi-continuous $\varphi\in D(\mathcal E)$.
Indeed, applying this to equilibrium potentials gives
\[
  \eta(O)\le C_\eta\,\operatorname{Cap}(O)^{1/2}
\]
for open $O$, hence $\eta$ charges no zero-capacity set.  Once
$\eta(N)=0$, for any $\pi$-null set $A$,
\[
  \eta(A)=\lambda\int_{N^c}R_\lambda(y,A)\,\eta(dy)=0,
\]
so $\eta\ll\pi$ and the $L^2(\pi)$ spectral gap forces $\eta=\pi$.
A possible way to prove the finite-energy inequality would be a
stationary Krylov/capacity estimate for the regularized processes with
constants depending on the bounded primitives $B_n$ rather than on
$\|G_{1/n}\|_\infty$.

Local-time form of the drift.  Bounebache--Zambotti write the SPDE with
a local-time-in-space term
\[
 -\frac12\int_\mathbb R f(da)\,\partial_\theta \ell^a_{t,\theta}.
\]
For $f=-2H_+$, $df=-2\delta_0$, so this term is
$\partial_\theta\ell^0_{t,\theta}$, which is formally
$\delta_0(u(t,\theta))$ by the occupation density formula in the spatial
variable.  Thus their Dirichlet form matches the Gibbs measure
$\pi\propto\exp(2\int_0^1{\bf1}_{\{v>0\}})\,d\gamma$, but their Markov
solution is constructed only outside a properly exceptional set.

Finite-dimensional projection/entropy route.  For an invariant law
$\eta$, finite-dimensional projections onto the first $m$ Dirichlet
modes should have Lebesgue densities because the projected dynamics has
a nondegenerate Brownian martingale part.  This by itself is insufficient
for $\eta\ll\pi$ in infinite dimensions.  A useful target would be
uniform relative entropy
\[
  \sup_m H((\Pi_m)_\#\eta\mid(\Pi_m)_\#\pi)<\infty,
\]
which would imply $\eta\ll\pi$ by projective convergence of entropy.
The difficulty is to derive such a bound from the stationary martingale
problem with the distributional drift; the drift contribution is a
surface/local-time measure rather than an $L^2$ vector field.

Roughness decomposition idea.  At fixed $t>0$, the linear stochastic
convolution has a Gaussian law equivalent to Brownian bridge.  If the
drift convolution were an independent Cameron--Martin shift, absolute
continuity would follow immediately.  In reality it is adapted and
depends on the same noise.  A possible approach is to split the driving
noise into independent high-frequency or late-time pieces and show that
the singular drift contribution cannot cancel the nondegenerate Gaussian
roughness.  This would amount to a Malliavin or conditional-density
argument for the singular/local-time SPDE.

\end{document}

